% ============================================================
%  Глава 1: Математические основы
% ============================================================

\section{Математические основы}
\label{sec:math}

\subsection{Рынок опционов и подразумеваемая волатильность}

Рассмотрим финансовый рынок с вероятностным пространством
$(\Omega, \FF, (\FF_t)_{t\geq 0}, \QQ)$, где $\QQ$~--- мартингальная мера.
Цена базового актива (например, S\&P~500) описывается процессом $(S_t)_{t\geq 0}$.

\begin{definition}[Европейский колл-опцион]
  Европейский колл-опцион с ценой исполнения $K$ и сроком погашения $T$ даёт
  держателю право (но не обязательство) купить базовый актив за $K$ в момент $T$.
  Его справедливая цена в момент $t=0$ при нулевой безрисковой ставке:
  \[
    C(S_0, K, T) = \EE^{\QQ}\!\bigl[(S_T - K)^+\bigr].
  \]
\end{definition}

\begin{definition}[Подразумеваемая волатильность]
  Подразумеваемая волатильность $\IV(K, T)$ --- это такое $\sigma > 0$, при
  котором формула Блэка--Шоулса воспроизводит наблюдаемую цену:
  \[
    C_{\BS}(S_0, K, T, \sigma) = C^{\mathrm{market}}(K, T).
  \]
  Как функция $(K, T)$, $\IV$ образует \textit{поверхность подразумеваемой
  волатильности}.
\end{definition}

Поверхность IV нетривиальна: для каждой даты погашения $T$ она образует
«улыбку» или «ухмылку» (volatility smile/skew) по страйку $K$.
Модель Блэка--Шоулса предсказывает плоскую поверхность, что противоречит
рыночным наблюдениям.

\subsection{Дробное броуновское движение и грубая волатильность}

\begin{definition}[Дробное броуновское движение]
  \textit{Дробное броуновское движение} (fractional Brownian motion, fBm)
  $B^H = (B^H_t)_{t\geq 0}$ с показателем Хёрста $H \in (0,1)$~--- это
  центрированный гауссовский процесс с ковариацией:
  \[
    \EE[B^H_s B^H_t]
    = \tfrac{1}{2}\!\left(s^{2H} + t^{2H} - |t-s|^{2H}\right).
  \]
  При $H = 1/2$ это стандартное броуновское движение; при $H < 1/2$~---
  \textit{грубый} (rough) процесс с траекториями более нерегулярными, чем BM.
\end{definition}

Empirically, Gatheral, Jaisson \& Rosenbaum (2018)~\cite{Gatheral2018}
показали, что показатель Хёрста реализованной волатильности~S\&P~500
оценивается как $\hat{H} \approx 0{,}08$~--- значительно ниже $1/2$,
что мотивирует использование $H < 1/2$ в моделях.

\begin{remark}
  При $H < 1/2$ фБД не является ни марковским, ни полумартингалом.
  Стохастическое исчисление для него требует специального обращения
  (исчисление Молчанова--Нуалара или репрезентация через ядро Вольтерра).
\end{remark}

\subsection{Модель Rough Heston}

\begin{definition}[Rough Heston~\cite{ElEuch2019}]
  В риск-нейтральной мере цена актива $S_t$ и мгновенная дисперсия $V_t$
  удовлетворяют:
  \begin{align}
    \frac{\dd S_t}{S_t}
      &= \sqrt{V_t}\,\dd W_t^S, \label{eq:S_dyn}\\
    V_t
      &= V_0 + \frac{1}{\Gamma(H+\tfrac{1}{2})}
         \int_0^t (t-s)^{H-1/2}\,\kappa(\theta - V_s)\,\dd s
         + \frac{\sigma}{\Gamma(H+\tfrac{1}{2})}
           \int_0^t (t-s)^{H-1/2}\,\sqrt{V_s}\,\dd W_t^V,
         \label{eq:V_dyn}
  \end{align}
  где $W^S$ и $W^V$~--- броуновские движения с корреляцией
  $\dd\langle W^S, W^V\rangle_t = \rho\,\dd t$.
\end{definition}

\textbf{Параметры модели}: $\theta = (\kappa, \theta_{\infty}, \sigma, \rho, V_0, H)$,
где:
\begin{itemize}
  \item $\kappa > 0$~--- скорость возврата к среднему;
  \item $\theta_{\infty} > 0$~--- долгосрочный уровень дисперсии;
  \item $\sigma > 0$~--- волатильность волатильности (vol-of-vol);
  \item $\rho \in (-1, 0)$~--- корреляция (отрицательна для equity);
  \item $V_0 > 0$~--- начальная мгновенная дисперсия;
  \item $H \in (0, 1/2)$~--- показатель Хёрста.
\end{itemize}

При $H = 1/2$ модель вырождается в классическую модель Heston (1993)~\cite{Heston1993}.

\subsection{Проблема идентифицируемости}

Шесть параметров модели нельзя восстановить независимо из типичного набора
наблюдаемых опционных цен~--- система близка к \textit{нерегулярной} (ill-posed).
Это проявляется в малых собственных значениях матрицы Фишера (FIM, подробнее
в~Главе~4): число обусловленности $\kappa(\mathbf{I}_{6\times 6}) \sim 10^6$,
что делает калибровку крайне чувствительной к шуму.

Анализ чувствительности (Delta-IV, Appendix~A) показывает, что из 6 параметров
лишь \textbf{три} определяют форму поверхности IV на горизонтах $T \in [0{,}1; 2{,}0]$:
\begin{align}
  v_0 &\equiv V_0 \qquad &&\text{(текущая дисперсия)}, \\
  \zeta &\equiv \sigma\rho \qquad &&\text{(наклон улыбки)}, \\
  \lambda &\equiv \sigma\sqrt{1 - \rho^2} \qquad &&\text{(кривизна улыбки)}.
\end{align}
Параметры $\kappa, \theta_{\infty}, H$ фиксируются при значениях
$\kappa^* = 1{,}0$, $\theta^* = 0{,}08$, $H^* = 0{,}08$
на основе исторических оценок. В Главе~4 обоснована эта редукция через FIM:
переход $6D \to 3D$ снижает число обусловленности в $1301\times$.

\subsection{Лифтинговое представление и дробное уравнение Риккати}

Ключевое преимущество Rough Heston~--- наличие \textbf{полуаналитической
характеристической функции}, доступной через решение дробного уравнения Риккати.
El~Euch~\&~Rosenbaum (2019) показали, что:
\[
  \EE^{\QQ}\!\left[e^{iu \ln S_T}\right]
  = \exp\!\left(V_0\,\Psi(u,T)
    + \kappa\theta_{\infty} \int_0^T \Psi(u,t)\,\dd t\right),
\]
где $\Psi(u, \cdot)$ удовлетворяет дробному уравнению Риккати:
\begin{equation}
  \label{eq:riccati_frac}
  \Psi(u, t)
  = \frac{iu(iu-1)}{2} + \int_0^t K(t-s)
    \left(\kappa - \sigma\rho\,iu\right)\Psi(u,s)\,\dd s
  - \frac{\sigma^2}{2}\int_0^t K(t-s)\,\Psi(u,s)^2\,\dd s,
\end{equation}
где ядро $K(t) = t^{H-1/2}/\Gamma(H+1/2)$.

\subsubsection{Приближение Бернштейна (лифтинг)}

Уравнение~\eqref{eq:riccati_frac} интегро-дифференциальное и не имеет
замкнутого решения. Approксимируем ядро $K$ суммой экспонент
(приближение Бернштейна, El~Euch~\&~Rosenbaum, 2019):
\[
  K(t) \approx K^N(t) = \sum_{i=1}^N c_i\,e^{-x_i t}, \qquad
  x_i = r_N^{i-1-N/2}, \quad
  c_i = \frac{x_i^{-(H+1/2)}}{\sum_{j} x_j^{-(H+1/2)}},
\]
где $r_N = 1 + 10 N^{-0{,}9}$. При $N = 40$ погрешность ядра:
\[
  \|K - K^{40}\|_{L^2[0,T]} \leq \varepsilon \approx 10\,\text{б.п. (IV)}.
\]
Замена ядра превращает уравнение~\eqref{eq:riccati_frac} в систему $N$ ОДУ
(\textit{поднятая} модель, lifted Heston):
\begin{equation}
  \label{eq:lifted}
  \frac{\dd \psi_i}{\dd t}
  = \frac{iu(iu-1)}{2}\,c_i
    - (\kappa + x_i)\psi_i
    - \sigma\rho\,iu \sum_j \psi_j
    - \frac{\sigma^2}{2}\!\left(\sum_j \psi_j\right)^{\!2}\!c_i, \quad i=1,\ldots,N.
\end{equation}
Характеристическая функция:
\[
  \ln\varphi(u,T) = V_0\!\sum_i \psi_i(T) + \kappa\theta_{\infty}\int_0^T\!\sum_i\psi_i(t)\,\dd t.
\]

\subsection{Условия отсутствия арбитража для поверхностей IV}

Не всякая поверхность $\IV(K,T)$ соответствует отсутствию арбитража.
Gatheral~\&~Jacquier (2011)~\cite{GatheralJacquier2011} дали необходимые и
достаточные условия в терминах тотальной дисперсии $w(k) = \IV^2(k,T)\cdot T$:

\begin{theorem}[Условия ГДЖ]
  Поверхность IV свободна от статического арбитража тогда и только тогда,
  когда выполнены:
  \begin{enumerate}
    \item \textbf{Calendar spread}: $w(k,T_1) \leq w(k,T_2)$ при $T_1 < T_2$
          (тотальная дисперсия не убывает по сроку).
    \item \textbf{Butterfly}: $g(k) \geq 0$, где
          \[
            g(k) = \left(1 - \frac{k\,w'(k)}{2w(k)}\right)^{\!2}
                   - \frac{(w'(k))^2}{4}\!\left(\frac{1}{w(k)} + \frac{1}{4}\right)
                   + \frac{w''(k)}{2}.
          \]
  \end{enumerate}
\end{theorem}

Эти условия используются как \textbf{регуляризующий штраф} при обучении FNO
(подробнее в~Главе~3, уравнение~\eqref{eq:arb_loss}).
